\newcommand{\id}{66}
\newcommand{\nom}{Fonctions combinatoires}
\newcommand{\sequence}{07}
\newcommand{\num}{01}
\newcommand{\type}{TP}
\newcommand{\descrip}{Conception d'un algorithme de calcul d'addition pour une calculatrice}
\newcommand{\competences}{A3-04: Identifier et décrire les liens entre les chaines fonctionnelles. \\ &  A3-07: Analyser un algorithme.  \\ &  A3-11: Interpréter tout ou partie de l'évolution temporelle d'un système séquentiel. \\ &  B2-12: Proposer une modélisation des liaisons avec leurs caractéristiques géométriques. \\ &  B2-13: Proposer un modèle cinématique paramétré à partir d'un système réel, d'une maquette numérique ou d'u \\ &  B2-14: Modéliser la cinématique d'un ensemble de solides. \\ &  B2-16: Modéliser une action mécanique. \\ &  B2-17: Simplifier un modèle de mécanisme. \\ &  B2-18: Modifier un modèle pour le rendre isostatique. \\ &  B2-20: Décrire le comportement d'un système séquentiel. \\ &  E2-01: Choisir un outil de communication adapté à l'interlocuteur.}
\newcommand{\nbcomp}{11}
\newcommand{\systemes}{Calculatrice de bureau}
\newcommand{\systemessansaccent}{Calculatrice de bureau}
\newcommand{\ilot}{1}
\newcommand{\ilotstr}{01}
\newcommand{\dossierilot}{\detokenize{Ilot_01 Calculatrice de bureau}}
\newcommand{\imageun}{Calculatrice_de_bureau}

\newcommand{\urlsysteme}{\href{https://www.costadoat.fr/systeme/85}{Ressources système}}
